\documentclass[10pt,a4paper]{amsart}
\usepackage{amsmath,amssymb,amsthm,graphicx}
\usepackage{verbatim}
\newenvironment{reference}{\comment}{\endcomment}
%\newenvironment{reference}{}{}
\newenvironment{slogan}{\comment}{\endcomment}
\newenvironment{history}{\comment}{\endcomment}

% For commutative diagrams we use Xy-pic
\usepackage[all]{xy}

% We use 2cell for 2-commutative diagrams.
\xyoption{2cell}
% Theorem environments.
%
\theoremstyle{plain}
\newtheorem{theorem}[subsection]{Theorem}
\newtheorem{proposition}[subsection]{Proposition}
\newtheorem{lemma}[subsection]{Lemma}

\theoremstyle{definition}
\newtheorem{definition}[subsection]{Definition}
\newtheorem{example}[subsection]{Example}
\newtheorem{exercise}[subsection]{Exercise}
\newtheorem{situation}[subsection]{Situation}

\theoremstyle{remark}
\newtheorem{remark}[subsection]{Remark}
\newtheorem{remarks}[subsection]{Remarks}

\numberwithin{equation}{subsection}

% Macros
%
\def\lim{\mathop{\mathrm{lim}}\nolimits}
\def\colim{\mathop{\mathrm{colim}}\nolimits}
\def\Spec{\mathop{\mathrm{Spec}}}
\def\Hom{\mathop{\mathrm{Hom}}\nolimits}
\def\Ext{\mathop{\mathrm{Ext}}\nolimits}
\def\SheafHom{\mathop{\mathcal{H}\!\mathit{om}}\nolimits}
\def\SheafExt{\mathop{\mathcal{E}\!\mathit{xt}}\nolimits}
\def\Sch{\mathit{Sch}}
\def\Mor{\mathop{\mathrm{Mor}}\nolimits}
\def\Ob{\mathop{\mathrm{Ob}}\nolimits}
\def\Sh{\mathop{\mathit{Sh}}\nolimits}
\def\NL{\mathop{N\!L}\nolimits}
\def\CH{\mathop{\mathrm{CH}}\nolimits}
\def\proetale{{pro\text{-}\acute{e}tale}}
\def\etale{{\acute{e}tale}}
\def\QCoh{\mathit{QCoh}}
\def\Ker{\mathop{\mathrm{Ker}}}
\def\Im{\mathop{\mathrm{Im}}}
\def\Coker{\mathop{\mathrm{Coker}}}
\def\Coim{\mathop{\mathrm{Coim}}}

% Boxtimes
%
\DeclareMathSymbol{\boxtimes}{\mathbin}{AMSa}{"02}

%
% Macros for moduli stacks/spaces
%
\def\QCohstack{\mathcal{QC}\!\mathit{oh}}
\def\Cohstack{\mathcal{C}\!\mathit{oh}}
\def\Spacesstack{\mathcal{S}\!\mathit{paces}}
\def\Quotfunctor{\mathrm{Quot}}
\def\Hilbfunctor{\mathrm{Hilb}}
\def\Curvesstack{\mathcal{C}\!\mathit{urves}}
\def\Polarizedstack{\mathcal{P}\!\mathit{olarized}}
\def\Complexesstack{\mathcal{C}\!\mathit{omplexes}}
% \Pic is the operator that assigns to X its picard group, usage \Pic(X)
% \Picardstack_{X/B} denotes the Picard stack of X over B
% \Picardfunctor_{X/B} denotes the Picard functor of X over B
\def\Pic{\mathop{\mathrm{Pic}}\nolimits}
\def\Picardstack{\mathcal{P}\!\mathit{ic}}
\def\Picardfunctor{\mathrm{Pic}}
\def\Deformationcategory{\mathcal{D}\!\mathit{ef}}
\usepackage{cleveref}
\crefname{originalreference}{original reference}{original references}
\Crefname{originalreference}{Original reference}{Original references}
\makeatletter
\newlabel{lemma-id-cone-null}{{09QK}{source}{}{sourcepointer-1}{}}
\newlabel{lemma-id-cone-null@cref}{{[originalreference][0][]09QK}{[1][0][]source}{}{}{}}
\makeatother
\crefname{originalreference}{original reference}{original references}
\Crefname{originalreference}{Original reference}{Original references}
\makeatletter
\newlabel{situation-ABC}{{09QJ}{source}{}{sourcepointer-2}{}}
\newlabel{situation-ABC@cref}{{[originalreference][0][]09QJ}{[1][0][]source}{}{}{}}
\makeatother
\title{Every morphism factors through an admissible monomorphism with a homotopy retraction}
\author{The Stacks Project Authors}
\date{stacks-project:dga.tex:local-snapshot-20261012; GNU FDL 1.2 or later}
\begin{document}
\maketitle
\noindent\textit{Editorial note.} All delivered mathematical text is the author's own, copied from the pinned Stacks Project chapter file sources/stacks-dga/pinned-original-source.tex; no statement or proof step was rewritten, reordered or removed.\par\medskip
\noindent\textit{Editorial note.} The counted proof is the author's complete argument: the direct-sum factorization object with its splittings, the assembly of tilde-alpha from the cone sequence, the admissible monomorphism verification with pi tilde-alpha = alpha and pi s = 1\_y, and the explicit homotopy computation via lemma-id-cone-null are all retained verbatim.\par\medskip
\noindent\textit{Editorial note.} The label lemma-factor names the counted statement, and the proof environment immediately following it in the source is delivered verbatim as its proof.\par\medskip
\noindent\textit{Editorial note.} The counted statement is carried by the author's own lemma environment in the chapter source.\par\medskip
\noindent\textit{Allowed standard prerequisites (editorial).}\begin{itemize}
\item Axiom (A): direct sums in the differential graded category. Same chapter, Situation situation-ABC, axiom (A) and Lemma lemma-id-cone-null. Provides the differential graded direct sum y plus C(1\_x) with its split projections and sections, and the lemma identifying the cone of the identity with a null object whose structure maps give the homotopy input.
\item The cone sequence of the identity. Same chapter, Definition definition-cone and the admissible short exact sequence x to C(1\_x) to x[1]. Supplies the natural morphism x to C(1\_x) fitting the admissible sequence, out of which the admissible monomorphism tilde-alpha is assembled.
\end{itemize}\medskip
\ifdefined\hypertarget\hypertarget{sourcepointer-1}{}\fi
\noindent\textit{Source reference.} Original reference 09QK: sources/stacks-dga/pinned-original-source.tex, source lines 3867--3867. Prerequisite pointer invoked at proof line(s) 4011; it names a standard fact of the same chapter that the counted argument uses, and every step of the counted proof built on it is retained verbatim in the delivered body.\par\medskip
\ifdefined\hypertarget\hypertarget{sourcepointer-2}{}\fi
\noindent\textit{Source reference.} Original reference 09QJ: sources/stacks-dga/pinned-original-source.tex, source lines 3762--3762. Prerequisite pointer invoked at proof line(s) 3939; it names a standard fact of the same chapter that the counted argument uses, and every step of the counted proof built on it is retained verbatim in the delivered body.\par\medskip
\begin{lemma}
\label{lemma-factor}
In Situation \ref{situation-ABC} let $\alpha : x \to y$
be a morphism in $\text{Comp}(\mathcal{A})$. Then there exists
a factorization in $\text{Comp}(\mathcal{A})$:
$$
\xymatrix{
x \ar[r]^{\tilde{\alpha}}  &
\tilde{y} \ar@<0.5ex>[r]^{\pi} &
y\ar@<0.5ex>[l]^s
}
$$
such that
\begin{enumerate}
\item $\tilde{\alpha}$ is an admissible monomorphism, and
$\pi\tilde{\alpha}=\alpha$.
\item There exists a morphism
$s:y\to\tilde{y}$ in $\text{Comp}(\mathcal{A})$
such that $\pi s=1_y$ and $s\pi$ is homotopic to $1_{\tilde{y}}$. 
\end{enumerate}
\end{lemma}

\begin{proof}
By axiom (A), we may let $\tilde{y}$ be the differential graded direct
sum of $y$ and $C(1_x)$, i.e., there exists a diagram
$$
\xymatrix@C=3pc{
y \ar@<0.5ex>[r]^s  &
y\oplus C(1_x) \ar@<0.5ex>[l]^{\pi} \ar@<0.5ex>[r]^{p} &
C(1_x)\ar@<0.5ex>[l]^t
}
$$
where all morphisms are of degree zero, and in
$\text{Comp}(\mathcal{A})$. Let $\tilde{y} = y \oplus C(1_x)$.
Then $1_{\tilde{y}} = s\pi + tp$. Consider now the diagram
$$
\xymatrix{
x \ar[r]^{\tilde{\alpha}}  &
\tilde{y} \ar@<0.5ex>[r]^{\pi} &
y\ar@<0.5ex>[l]^s
}
$$
where $\tilde{\alpha}$ is induced by the morphism $x\xrightarrow{\alpha}y$
and the natural morphism $x\to C(1_x)$ fitting in the admissible
short exact sequence
$$
\xymatrix{
x \ar@<0.5ex>[r]  &
C(1_x) \ar@<0.5ex>[l] \ar@<0.5ex>[r] &
x[1]\ar@<0.5ex>[l]
}
$$
So the morphism $C(1_x)\to x$ of degree 0 in this diagram,
together with the zero morphism $y\to x$, induces a degree-0
morphism $\beta : \tilde{y} \to x$. Then $\tilde{\alpha}$ is an
admissible monomorphism since it fits into the admissible short
exact sequence
$$
\xymatrix{
x\ar[r]^{\tilde{\alpha}} &
\tilde{y} \ar[r] &
x[1]
}
$$
Furthermore, $\pi\tilde{\alpha} = \alpha$ by the construction of
$\tilde{\alpha}$, and $\pi s = 1_y$ by the first diagram. It
remains to show that $s\pi$ is homotopic to $1_{\tilde{y}}$.
Write $1_x$ as $d(h)$ for some degree $-1$ map. Then, our
last statement follows from
\begin{align*}
1_{\tilde{y}} - s\pi
= &
tp \\
= &
t(dh)p\quad\text{(by Lemma \ref{lemma-id-cone-null})} \\
= &
d(thp)
\end{align*}
since $dt = dp = 0$, and $t$ is of degree zero. 
\end{proof}

\end{document}
